\section{Theoretical Energy Consumption}

\subsection*{Predicted Energy Costs Across Multiple Platforms}

Using the \gls{flop} and \gls{sop} count recorded in Table \ref{tab:layer_ops} with the energy constants from Table \ref{tab:energy_constants}, we can predict the theoretical energy consumption of the \texttt{DetectorNX-G3-SNN} architecture across multiple platforms:

\begin{table}[htbp]
\centering
\begin{tabular}{|l|c|c|}
\hline
\textbf{Hardware Platform} & \textbf{Energy Per Frame ($mJ$)} & \textbf{Energy Savings (\%)} \\
\hline
NVIDIA A100 (\gls{cnn} Baseline) & 180.8576 & --- \\
Intel Loihi 2 (\gls{snn}) & 43.0223 & 76.21\% \\
IBM TrueNorth (\gls{snn}) & 69.4912 & 61.58\% \\
\hline
\end{tabular}
\caption[Neuromorphic Hardware Energy Extrapolation]{Extrapolated energy consumption per frame across different hardware architectures. Savings are calculated relative to the NVIDIA A100 GPU baseline.}
\label{tab:hardware_extrapolation}
\end{table}

\subsection*{Critical Analysis}

Deploying the \gls{snn} on a next-generation Intel Loihi 2 processor yields a projected $76.21\%$ reduction in energy consumption per frame compared to running the continuous baseline on the NVIDIA A100 GPU used for training. This outcome validates the hypothesis formulated in Section \ref{subsubsec:energy_hypothesis}.

This energy gain is realised by two compounding factors. 
First, the biological sparsity of the network results in a massive reduction in computational workload: the \gls{snn} requires only $2.67$ G-\glspl{sop} per frame, compared to the $9.22$ G-\glspl{flop} required by the \gls{cnn} baseline----a $71.02\%$ reduction in total operations. 
Second, the advanced architecture of Loihi 2 reduces the energy cost of each remaining operation to just $16.1$\,pJ \cite{gralewicz2024spikes}, eliminating the first-generation 'routing overhead' penalty that previously constrained neuromorphic hardware.

Furthermore, this 76.21\% energy savings was achieved despite the integration of the high-resolution G3 backbone ($22 \times 40$ feature maps). This proves that the biological sparsity of the \gls{snn} is robust enough to sustain high-definition spatial reasoning without inducing a catastrophic power surge---a scaling feat that is mathematically impossible for dense \gls{cnn} architectures.

Finally, it must be noted that comparing experimental neuromorphic silicon against a mature, massively parallel and optimised NVIDIA \gls{gpu} gives us an upper-bound on the energy savings of event-driven \glspl{snn}, rather than a realistic value;
deploying models such as \texttt{DetectorNX} on real neuromorphic hardware requires substantial development on its own, and the power saving of $76.21\%$ on Loihi 2 is unlikely to be realised in the real-world due to unaccounted overheads.
As neuromorphic hardware develops, the per-\gls{sop} energy cost will continue to converge towards the theoretical \acrfull{cmos} limit established by Horowitz \cite{horowitz20141}, and the power savings will scale exponentially higher towards the optimal $94.33\%$ value obtained in Section \ref{subsubsec:mean_energy_savings}.

Overall, the predicted power consumption on the various platforms described in Table \ref{tab:hardware_extrapolation} serves as conclusive proof that event-driven \glspl{snn} can bring highly power-efficient computing to resource-constrained edge environments, validating the neuromorphic paradigm as a viable, sustainable solution for real-time autonomous traffic monitoring.

\subsection*{Energy Efficiency Comparison}
\label{subsec:efficiency_landscape}

The following experiment evaluates the energy efficiency frontier, defined as $\text{efficiency} = \frac{mAP_{50}}{\text{Power Per Frame}}$,
of \texttt{DetectorNX-G3-SNN} against contemporary published \gls{snn} detectors.

This comparison serves as an illustrative guide rather than a strict benchmark, as it is confounded by three fundamental asymmetries:
\begin{enumerate}
    \item \textbf{Task Complexity:} 
    Competitors evaluate on the 80-class COCO dataset, whereas \texttt{DetectorNX} performs single-class vehicle detection on \gls{etram}. 
    As a result, the baseline \gls{map_50} scores measure different joint-classification difficulties.
    \item \textbf{Input Resolution:} The COCO-based models typically operate at lower spatial dimensions (e.g., SpikeDet processes $640 \times 640$ inputs \cite{fan2025spikedet}). In contrast, \texttt{DetectorNX} processes high-definition automotive data at $1280 \times 720$, representing a $2.25\times$ larger raw pixel volume per frame.
    \item \textbf{Domain Statistics:} 
    COCO features dense RGB backgrounds, whereas \gls{etram} provides sparse, dynamic event-driven streams.
\end{enumerate}

\begin{table}[htbp]
\centering
\resizebox{\textwidth}{!}{%
\begin{tabular}{|l|l|c|c|c|c|c|}
\hline
\textbf{Model} & \textbf{Dataset (Task)} & \textbf{\gls{map_50}} & \textbf{Params (M)} & \textbf{G-SOPs} & \textbf{Power (mJ)} & \textbf{mAP/mJ$^\dagger$} \\
\hline
SpikeDet-S \cite{fan2025spikedet}    & COCO (80-class) & 62.6 & 22.0 & 20.3  & 19.6  & 3.19  \\
SpikeYOLO-N \cite{spike_yolo}        & COCO (80-class) & 59.2 & 13.2 & 25.1  & 23.1  & 2.56  \\
EMS-YOLO \cite{ems_yolo}             & COCO (80-class) & 50.1 & 33.9 & 29.2  & 29.0  & 1.73  \\
SpikeDet-X \cite{fan2025spikedet}    & COCO (80-class) & 67.6 & 75.2 & 42.9  & 40.4  & 1.67  \\
E-SpikeFormer \cite{espikeformer}    & COCO (80-class) & 58.8 & 38.7 & 127.0 & 119.5 & 0.49  \\
\hline
\multicolumn{7}{|c|}{\itshape Below: evaluated on a different task --- figures not directly comparable.} \\
\hline
\textbf{\texttt{DetectorNX-G3-SNN}}  & \textbf{ETraM (1-class)} & \textbf{59.8} & \textbf{11.7} & \textbf{2.7} & \textbf{2.4} & \textbf{24.92} \\
\hline
\end{tabular}%
}
\caption[Comparative efficiency landscape of related SNN detectors.]{Comparative efficiency landscape for \texttt{DetectorNX-G3-SNN} against published \gls{snn} object detectors. All competitor figures are reported on the COCO dataset (80-class), while \texttt{DetectorNX-G3-SNN} is evaluated on the single-class vehicle subset of \gls{etram}. G-SOPs reflects total synaptic operations (MAC + AC) per frame; Power reflects estimated energy per frame on neuromorphic hardware. Competitors are sorted by descending mAP/mJ ratio. $^\dagger$ The mAP/mJ ratio is not directly comparable across rows of differing dataset, since the numerator scales with task difficulty (80-class versus 1-class detection).}
\label{tab:efficiency_landscape}
\end{table}

Despite these asymmetries, the data in Table~\ref{tab:efficiency_landscape} reveals a compelling observation. At just $2.4$\,mJ per frame, \texttt{DetectorNX-G3-SNN} is approximately an order of magnitude more energy-efficient than the surveyed alternatives, whilst simultaneously carrying the smallest parameter budget ($11.7$\,M) and processing the highest-resolution input.

This extreme efficiency is driven by the G3 architecture. 
The \texttt{SEWResBlock} layout propagates spikes without inducing an activation explosion, while the \gls{frr} constraint strictly bounds layer-wise firing rates below $20\%$. Combined with a deliberately lightweight multi-box head versus the heavy transformer head of E-SpikeFormer \cite{espikeformer}, the model ensures the global \gls{sop} count remains minimal regardless of input scale.

The honest takeaway is not a direct claim of superior efficiency, given the dataset mismatches, but rather that \texttt{DetectorNX-G3-SNN} occupies a uniquely low-power, high-resolution region of the SNN landscape. The ability to maintain $\sim 60\%$ \gls{map_50} at HD resolutions for under $3$\,mJ suggests this sparsity-first architecture could scale highly effectively to multi-class tasks; this hypothesis is left for future COCO-based validation.